\section{Prospective hard-UC confirmation}\label{app:hard-confirm}
This separate synthetic study tests whether soft-UC supervision improves a prospectively selected hard-UC complete system over native membership heads under equal measured training/development allocation. Its protocol and two primary contrasts were frozen before this new run, without external preregistration. It was motivated by the post hoc diagnostic in Appendix~\ref{app:uc-equal-compute}; its fresh collections and inference family are not pooled with that study or the matched-epoch comparisons. The previous soft-set null remains unchanged.

\subsection{Data, training, and decoder-aligned selection}
Twelve independent collections each contain 512 training and 128 development graphs at $n=12$, and 256 tests at each of $n=24,48$. The five generator families, heterogeneous observation counts, planned missingness 0.25, and tie probability 0.20 match Appendix~\ref{app:uc-equal-compute}. There are 6,144 distinct held-out graphs and three paired initializations per collection. The pair architecture, Adam learning rate 0.001, batch size 16, gradient clipping 5, float32 execution, and canonical UC temperature 0.035 are unchanged. Paired initial states and minibatch-order seeds are shared; equal time permits different update counts. Candidate order rotates across systems.

Six systems are independently trained and timed: pairwise-only with hard UC; ordinary and relational auxiliary supervision with native heads; STE soft-UC supervision with hard UC; and separately trained ordinary/relational auxiliary controls with hard UC. Auxiliary controls retain their membership-head training loss. STE differentiates the original soft operator, not hard UC. All hard systems select through their own hard decoder, rather than reusing the native-head choices. This controls decoder-specific selection; structural-loss alignment remains part of the treatment.

Hard UC applies the unchanged covering oracle to each stored direction $\widehat P_{ab}>1/2$. There is no membership-score cutoff, per-test-size calibration, or oracle cardinality. Float32 reciprocal probabilities can round asymmetrically at $1/2$; if neither direction exceeds $1/2$, both arcs are omitted. The resulting graph can be incomplete. Upper-triangle exact ties, genuinely unoriented pairs, and one-sided reciprocal rounding are recorded separately.

Each supervised system fits two loss weights, $\lambda\in\{0.1,1\}$, for 110 seconds each; pairwise-only fits one candidate for 220 seconds. Five checkpoints target allocation fractions 0.025, 0.1, 0.25, 0.5, and 1. Selection maximizes development per-graph F1 on selective non-singleton UC at $n=12$. Native heads also optimize the original 43 absolute cutoffs. Exact ties prefer an earlier checkpoint, then smaller $\lambda$, and for heads a cutoff nearer 0.5 then smaller cutoff. Every system's three initialization choices are frozen and hashed before that collection's test graphs are generated.

\subsection{Budget and inference validity}
The study contains 396 candidate fits and 216 selected complete systems. Each system receives 220 charged seconds per collection/initialization, totaling 47,520 allocated seconds. Charged work includes synchronized training, native-decoder development prediction/selection, checkpoint serialization, in-loop bookkeeping, and final winner selection/freezing. Shared final serialization cost is apportioned equally across systems. Data generation, initialization/transfers, common warmup, post-budget integrity bookkeeping, and held-out evaluation are recorded separately. The checkpoint scheduler reserves the largest prior checkpoint cost at each synchronized minibatch boundary. All candidates and complete-system totals satisfy the prospectively fixed $\pm2\%$ gate; charged complete-system time totals 47,574.23 seconds. This satisfies the specified allocation tolerance, not total application runtime or active parameter count.

The two primary contrasts are STE-hard minus ordinary native-head F1 and minus relational native-head F1 on selective non-singleton UC at $n=24$. Eligible graphs average within each initialization, then the three initializations within collection. All twelve independent paired collection differences enter inference. Individual 95\% intervals use $t_{11}$ and approximate normality of collection differences; two-sided exhaustive sign flipping enumerates all 4,096 patterns and assumes null exchangeability under independent sign reversals. Holm adjustment covers exactly these two contrasts. Secondary sizes, hard controls, exact recovery, relation errors, and cardinalities are descriptive.

\begin{table}[ht]\centering\small
\begin{tabular*}{\linewidth}{@{\extracolsep{\fill}}lrrrr@{}}\toprule
Native comparator & Difference & 95\% CI & Raw $p$ & Holm $p$\\\midrule
Ordinary head & 0.03493 & $[0.02383,0.04602]$ & 0.0004883 & 0.0009766\\
Relational head & 0.03843 & $[0.02763,0.04922]$ & 0.0004883 & 0.0009766\\
\bottomrule\end{tabular*}

\caption{The separate frozen hard-UC primary family. Differences are STE-hard minus native heads at selective $n=24$. Both comparisons retain all twelve collection effects. Intervals are individual paired $t_{11}$ intervals; Holm covers the two exhaustive sign-flip tests.}\label{tab:hard-confirm-primary}
\end{table}

\begin{table}[ht]\centering\scriptsize
\begin{tabular*}{\linewidth}{@{\extracolsep{\fill}}llrrrrrr@{}}\toprule
$n$ & System / readout & All F1 & Selective F1 & Exact (\%) & Selected & Reference & Pair MSE\\\midrule
24 & Pair / hard & 0.5131 & 0.5865 & 1.45 & 16.72 & 9.57 & 0.0407\\
24 & Ordinary / native & 0.5902 & 0.6327 & 2.02 & 16.34 & 9.57 & 0.0445\\
24 & Relational / native & 0.5832 & 0.6292 & 2.53 & 16.43 & 9.57 & 0.0398\\
24 & STE / hard & 0.6130 & 0.6676 & 1.16 & 10.75 & 9.57 & 0.0354\\
24 & Ordinary / hard & 0.5231 & 0.5972 & 1.41 & 15.87 & 9.57 & 0.0386\\
24 & Relational / hard & 0.5621 & 0.6237 & 1.44 & 14.21 & 9.57 & 0.0402\\
\midrule
48 & Pair / hard & 0.4170 & 0.4012 & 0.36 & 40.31 & 13.13 & 0.0405\\
48 & Ordinary / native & 0.4725 & 0.4309 & 0.00 & 35.22 & 13.13 & 0.0434\\
48 & Relational / native & 0.4667 & 0.4318 & 0.48 & 35.17 & 13.13 & 0.0377\\
48 & STE / hard & 0.5272 & 0.5159 & 1.75 & 25.99 & 13.13 & 0.0345\\
48 & Ordinary / hard & 0.4293 & 0.4121 & 1.02 & 38.54 & 13.13 & 0.0378\\
48 & Relational / hard & 0.4554 & 0.4324 & 1.10 & 35.57 & 13.13 & 0.0378\\
\bottomrule\end{tabular*}

\caption{All six independently timed systems. All-class F1 uses the natural mixture; remaining metrics use selective non-singleton cores. Exact is complete-set recovery (\%). Pair MSE uses latent simulated truth unavailable to the predictor. Means average graphs, initializations, and twelve collections. $n=48$ and the common-hard controls are descriptive.}\label{tab:hard-confirm-systems}
\end{table}

STE-hard improves primary F1 by 0.03493/0.03843 over ordinary/relational native heads (Table~\ref{tab:hard-confirm-primary}); all twelve differences are positive for each contrast. The common-hard controls are also below STE-hard descriptively, and STE has lower latent pair MSE. These findings support relation learning and the tested complete hard-UC pipeline. They do not establish general architectural superiority or erase the earlier soft-set null. Exact selective recovery is only 1.16\% at $n=24$, below native-head recovery of 2.02/2.53\%. At $n=48$, STE-hard selects 25.99 alternatives against a 13.13-member reference, despite its higher F1. Successful overlap remains distinct from accurate membership and cardinality.

\subsection{Saved evidence and replay}
Production uses an NVIDIA L40S with PyTorch 2.6.0+cu124/CUDA 12.4, NumPy 1.26.4, SciPy 1.13.1, and pandas 2.2.3. TF32 and mixed precision are disabled. The archive retains observations, latent labels/probabilities, paired initial states, candidate histories and budgets, development predictions/choices, selected checkpoints, and test predictions. Selected-checkpoint CUDA replay verifies 432 prediction units containing 110,592 case outputs. An independent raw-output audit recomputes all 110,592 selected predictions, checks 216 development choices and 396 candidate/216 system timing records, and reproduces the primary statistics and all reported system summaries. Archive payload hashes and CRCs pass for both the full and review copies. Replayed stored states and independently reconstructed metrics support computational consistency, not independent retraining, optimizer convergence, formal execution attestation, or external replication.
